\documentclass[11pt]{article}
\usepackage{amsmath,amssymb,amsthm,hyperref}
\newtheorem{theorem}{Theorem}
\newtheorem{lemma}[theorem]{Lemma}
\newcommand{\E}{\mathbb E}
\newcommand{\Prob}{\mathbb P}
\newcommand{\Var}{\operatorname{Var}}
\title{A sharp finite-regime lower bound for random palindrome blocks}
\author{Research proof; two independent full reviews completed}
\date{30 September 2026}
\begin{document}
\maketitle

Write $R_\sigma=n!P(\sigma)$ for $m$ independent uniformly relabelled
palindrome blocks on $n$ labels. This note completes the lower bound for
this random family, including the logarithm, by comparing small assignment
problems with Gaussian ones. It concerns this family only.

\begin{theorem}\label{thm:main}
There are absolute constants $c,C>0$ and $n_0$ such that, for $n\geq n_0$
and
\[
 m\geq Cn^{13/10}(\log n)^{1/5},
\]
\[
 \Prob\left(\max_{\sigma\in S_n}|R_\sigma-1|
       \geq c\frac{n^{13/4}\sqrt{\log n}}{m^{5/2}}\right)
 \geq1-Cn^{-9/10}.
\]
Thus the previously proved sufficient scale
$n^{13/10}(\log n)^{1/5}\varepsilon^{-2/5}$ matches the typical necessary
scale for this construction in the large-block, small-error regime.
\end{theorem}

We use the notation and elementary estimates from the companion notes.
Set
\[
 d=n-3,\quad w_i(a)=\binom da p_i^a(1-p_i)^{d-a},\quad
 p_i=(i-1)/(m-1),\quad
 h=\binom n3m^{-3}(1-1/m)^d.
\]
The triple projection is
$S_\sigma=h\sum_i\sum_aw_i(a)\delta_{i,a}$, with
$\delta(l,x,r)=\tfrac12-\tfrac32\mathbf1\{x<\min(l,r)\}$.
For central starts $a\in[d/4,3d/4]$ and $m\geq n$,
\begin{equation}\label{eq:weights}
 \max_iw_i(a)\leq Cn^{-1/2},\qquad
 c m n^{-3/2}\leq\sum_iw_i(a)^2\leq C m n^{-3/2}.
\end{equation}
These bounds follow from Stirling, the beta integral, and the Riemann-sum
estimate for a unimodal function, as detailed in
\emph{A matching power lower bound for random palindrome blocks}.

\section{A smooth comparison inequality}
\begin{lemma}\label{lem:comparison}
Let $X_i$ be independent centered random vectors in $\mathbb R^D$, and
let independent centered Gaussian vectors $G_i$ have the respective same
covariance matrices. For $\beta>0$,
\[
 \left|\E F_\beta\Bigl(\sum_iX_i\Bigr)
       -\E F_\beta\Bigl(\sum_iG_i\Bigr)\right|
 \leq C\beta^2\sum_i\bigl(\E\|X_i\|_\infty^3+
                              \E\|G_i\|_\infty^3\bigr),
\]
where $F_\beta(x)=\beta^{-1}\log\sum_{\alpha=1}^D e^{\beta x_\alpha}$.
Also $\max_\alpha x_\alpha\leq F_\beta(x)
\leq\max_\alpha x_\alpha+\beta^{-1}\log D$.
\end{lemma}
\begin{proof}
For any direction $v$, the third directional derivative of $F_\beta$ is
$\beta^2$ times the third centered moment of the numbers $v_\alpha$ under
the softmax probabilities. Its absolute value is at most
$8\beta^2\|v\|_\infty^3$. Replace the vectors one at a time and expand
along the replaced vector through degree two. The first two Taylor terms
cancel because means and covariance matrices match. The stated third-order
remainder bound proves the lemma. The maximum inequalities are immediate.
\end{proof}

\section{One small assignment problem}
Put $L=\lfloor n^{1/10}\rfloor$. Inside a window of $3L+4$ consecutive
central positions, choose $L$ movable positions spaced by three, leaving
two positions of padding at each end. Keep the other labels fixed and
permute the $L$ movable labels arbitrarily among the movable positions.
There are $D=L!$ choices $\pi$. Exactly three triple starts around each
movable position can be affected. Their set $\mathcal A$ has size $3L$,
and all corresponding triples lie inside this window.

Let $T_\pi$ be the contribution of these starts to $S$ under assignment
$\pi$, and normalize by
\[
 s_0=h\sqrt{m/n^{3/2}},\qquad
 X_i(\pi)=\frac h{s_0}\sum_{a\in\mathcal A}
                 w_i(a)\delta_{i,a}^{\pi}.
\]
Then $T_\pi/s_0=\sum_iX_i(\pi)$, a sum of independent centered vectors.
Write $M_i=\max_{a\in\mathcal A}w_i(a)$. From \eqref{eq:weights},
\begin{equation}\label{eq:Mthird}
 \sum_iM_i^3\leq\sum_{a\in\mathcal A}\sum_iw_i(a)^3
 \leq CL\frac m{n^2}.
\end{equation}
The deterministic bound on a single vector is
\[
 \|X_i\|_\infty\leq C L\frac{n^{3/4}}{\sqrt m}M_i.
\]
For each fixed target, valley covariances vanish beyond distance two and
have uniformly bounded absolute row sums. Thus
\[
 \max_\pi\Var X_i(\pi)\leq CL\frac{n^{3/2}}m M_i^2.
\]
A Gaussian vector with coordinate variances at most $v^2$ satisfies
$\E\|G\|_\infty^3\leq Cv^3(\log(2D))^{3/2}$, by the Gaussian tail bound
and a union bound. Since $\log(2D)\leq CL\log L$, \eqref{eq:Mthird} gives
\begin{equation}\label{eq:third}
 \sum_i\bigl(\E\|X_i\|_\infty^3+\E\|G_i\|_\infty^3\bigr)
 \leq C L^4(\log L)^{3/2}\frac{n^{1/4}}{\sqrt m}.
\end{equation}

The Gaussian vector $G=\sum_iG_i$ contains an independent assignment
matrix. Specifically, let $u(l,x,r)$ be the canonical three-variable
Hoeffding component of $\delta$, for which $\E u^2=1/20$.
For each movable position $v$ and movable label $b$, project onto
\[
 \frac h{s_0}\bigl(w_i(v-1)
        u(Y_{i,\lambda_v},Y_{i,b},Y_{i,\rho_v})\bigr)_{i=1}^m,
\]
where $\lambda_v,\rho_v$ are the two fixed adjacent anchors. Distinct such
kernels are orthogonal, since their triples share at most two variables
and $u$ is centered given any two. The corresponding Gaussian entries
$Z_{v,b}$ are independent, and their variances satisfy
\[
 \Var Z_{v,b}=\frac{h^2}{20s_0^2}\sum_iw_i(v-1)^2\asymp1
\]
by \eqref{eq:weights}. Exactly as in the companion Gaussian-limit note,
\[
 G(\pi)=\sum_vZ_{v,\pi(v)}+H(\pi),
\]
with the centered Gaussian residual independent of the matrix. No limiting
argument is used here: this is an orthogonal decomposition of the finite
Gaussian surrogate with the given covariance matrix.
Conditional Jensen and a greedy row-by-row assignment therefore give
\begin{equation}\label{eq:gaussianmax}
 \E\max_\pi G(\pi)\geq cL\sqrt{\log L}.
\end{equation}
Indeed the first $L/2$ rows each have at least $L/2$ fresh independent
Gaussian columns with standard deviations bounded below, and later greedy
increments have nonnegative expectations.

Apply Lemma~\ref{lem:comparison} with
$\beta=B\sqrt{\log L}$, where a sufficiently large absolute $B$ makes
the smoothing error $(\log D)/\beta$ less than one fourth of
\eqref{eq:gaussianmax}. The comparison error, divided by
$L\sqrt{\log L}$, is at most
\[
 C L^3(\log L)^2\frac{n^{1/4}}{\sqrt m}
 \leq C n^{-1/10}(\log n)^2=o(1)
\]
whenever $m\geq n^{13/10}$. Therefore, uniformly in this range,
\begin{equation}\label{eq:chunkmean}
 \E\bigl(\max_\pi T_\pi-\min_\pi T_\pi\bigr)
 \geq c s_0L\sqrt{\log L}.
\end{equation}
We used only $\E\min_\pi T_\pi\leq0$, since each coordinate is centered.

We also need a second-moment upper bound. For any fixed $\pi$,
\[
 \sum_i\Var X_i(\pi)\leq CL,\qquad
 |X_i(\pi)|\leq C L n^{1/4}/\sqrt m.
\]
The elementary Bernstein moment estimate for a sum of bounded independent
centered variables gives, for $q\geq2$,
\[
 \left\|\sum_iX_i(\pi)\right\|_q
 \leq C\left(\sqrt{qL}+qL n^{1/4}/\sqrt m\right).
\]
This estimate follows by expanding each exponential moment, summing the
geometric remainder using boundedness, and integrating the resulting
Bernstein tail bound. Take $q=\max(2,\log(2D))$ and use
$\|\max_\pi|X(\pi)|\|_q\leq D^{1/q}\max_\pi\|X(\pi)\|_q$.
Since $L\leq n^{1/10}$ and $m\geq n^{13/10}$, this yields
\begin{equation}\label{eq:chunksecond}
 \E\bigl(\max_\pi T_\pi-\min_\pi T_\pi\bigr)^2
 \leq C s_0^2L^2\log L.
\end{equation}

\section{Independent windows and the full distribution}
Fit $Q\asymp n/L$ disjoint windows of the above type in the middle third
of a base target. Permute their movable labels independently. The family
$\mathcal F$ has $(L!)^Q$ members, at most $e^{n\log n}$.
Every affected triple lies wholly in one window, so the range of the
triple projection over this family is exactly the sum of the $Q$ individual
ranges. These ranges are independent because their label sets are disjoint
in every block. From \eqref{eq:chunkmean}, \eqref{eq:chunksecond}, and
Chebyshev, with probability $1-O(Q^{-1})$,
\[
 \max_{\sigma\in\mathcal F}|S_\sigma|
 \geq c Qs_0L\sqrt{\log L}
 \geq c\frac{n^{13/4}\sqrt{\log n}}{m^{5/2}}.
\]
Here $h\asymp n^3/m^3$, $\log L\asymp\log n$, and
$Q^{-1}=O(n^{-9/10})$.

Put $A=n^{13/4}\sqrt{\log n}/m^{5/2}$ and $t=m/n$.
The coefficient estimates in the refined upper-bound note, with
$p=4n\log n$, give for $U_\sigma=R_\sigma-1-S_\sigma$
\[
 \|U_\sigma\|_p\leq C_1 A(A+t^{-1})
\]
for sufficiently large $n$ and $m\geq C_0n^{13/10}(\log n)^{1/5}$.
For completeness, short-support levels $\ell\geq2$ contribute
$O(A^2)$; the short-support first level of lengths $k\geq4$ contributes
$O(A/t)$; the long-support tilted estimate is
$\sqrt2\exp(-n\log t/12+4pn/t^4)\leq A/t$.
These are the same coefficient-level estimates used in the earlier
power-only lower bound, with the larger value of $p$.

A union bound over $\mathcal F$ and Markov at threshold
$e C_1A(A+t^{-1})$ fail with probability at most $e^{-3n\log n}$.
Taking the absolute threshold constant $C_0$ large enough makes that
remainder at most half of the triple-projection lower bound, because
$A\leq C_0^{-5/2}$ and $t^{-1}\leq C_0^{-1}$.
This proves Theorem~\ref{thm:main}.

\paragraph{A threshold consequence.}
For some fixed absolute $\varepsilon_0>0$, if
$m=o(n^{13/10}(\log n)^{1/5})$, the probability of relative
$\varepsilon_0$-fairness tends to zero. To see this, when $m$ is at least
linear in $n$, restrict to $k$ labels with
$m\asymp Ck^{13/10}(\log k)^{1/5}$ and apply the theorem; then
$k\to\infty$, $k\leq n$, and the lower error is an absolute constant.
For $m<n/(4e)$, outcome counting already forces an absent permutation.
The remaining cases have $m\to\infty$ and admit the same restriction.
Relative fairness is preserved by restriction because marginal relative
errors are averages of full-order relative errors.

\paragraph{Dependencies and scope.}
The elementary binomial estimates and full nonlinear remainder estimates
are proved in the companion notes
\emph{A refined uniform bound for random palindrome blocks} and
\emph{A matching power lower bound for random palindrome blocks}.
The canonical variance $1/20$ and Gaussian assignment decomposition are
proved in \emph{The logarithmic factor in the large-block Gaussian limit}.
Here that decomposition is applied to a finite covariance-matched Gaussian
surrogate, and Lemma~\ref{lem:comparison} supplies the quantitative step.
No lower bound for general approximate dice constructions follows.
\end{document}
